\section{Metodología}

\subsection{Las tres preguntas}

\begin{table}[htbp]
\centering
\caption{Preguntas analíticas del proyecto}
\label{tab:preguntas}
\begin{tabularx}{\textwidth}{clX}
\toprule
\# & Pregunta & Salida \\
\midrule
Q1 & ¿Cómo varía la demanda de viajes por hora del día y día de la semana? & 168 combinaciones día $\times$ hora \\
Q2 & ¿Qué estaciones tienen mayor desequilibrio entre llegadas y salidas? & 808 estaciones con actividad \\
Q3 & ¿Existen rutas origen-destino que concentran gran parte del tráfico? & 510,672 rutas distintas \\
\bottomrule
\end{tabularx}
\end{table}

\subsection{Las tres herramientas}

\begin{table}[htbp]
\centering
\caption{Herramientas comparadas}
\label{tab:herramientas}
\begin{tabularx}{\textwidth}{lXX}
\toprule
Herramienta & Arquitectura & Por qué está en la comparación \\
\midrule
Pandas & DataFrame en memoria, columnar sobre NumPy, mono-hilo & Caso base: rápido mientras el dataset quepa en RAM \\
DuckDB & Motor OLAP columnar embebido, vectorizado y multi-hilo & Diseñado exactamente para este tipo de agregación \\
SQLite & Motor OLTP embebido, orientado a filas & Misma interfaz SQL, arquitectura de almacenamiento opuesta \\
\bottomrule
\end{tabularx}
\end{table}

SQLite es la tercera herramienta y aporta el contraste más informativo del trabajo: comparte el lenguaje con DuckDB pero no su modelo de almacenamiento, así que aísla el efecto de la arquitectura frente al del lenguaje de consulta.

\subsection{Protocolo de ejecución independiente}

Cada combinación de versión, herramienta y pregunta es una \textbf{ejecución independiente que arranca desde el archivo}: se carga desde cero y después se corre la consulta. El tiempo reportado es carga $+$ query, y ese es el número que se compara contra el límite de 10 minutos.

La razón es que SQLite no conserva el archivo cargado entre consultas, a diferencia de Pandas, que mantiene el DataFrame en RAM, y de DuckDB, que mantiene la tabla en memoria. Aplicar el mismo protocolo a las tres evita que alguna se beneficie de reutilizar una carga que otra tuvo que repetir.

Otras reglas que sostienen la honestidad de la medición:

\begin{itemize}
  \item La carga a SQLite \textbf{no usa Pandas}. Se lee el CSV en streaming con el módulo \texttt{csv} de la librería estándar y se inserta por lotes de 200,000 filas con \texttt{executemany}. Usar \texttt{df.to\_sql()} mediría a Pandas disfrazado de SQLite.
  \item DuckDB \textbf{materializa} la tabla con \texttt{CREATE TABLE ... AS SELECT}, en vez de leer el CSV de forma perezosa, para competir en igualdad de condiciones con \texttt{pd.read\_csv()}.
  \item La base SQLite se crea desde cero en cada ejecución y se borra al terminar.
  \item El pico de RAM se captura con un hilo que muestrea el proceso cada 50 ms. Medir solo antes y después perdería el pico transitorio que ocurre dentro de la lectura.
  \item Los días de la semana se normalizan a lunes $=0$ en las tres herramientas, porque cada motor los numera distinto.
\end{itemize}

\subsection{Validación cruzada}

Un benchmark solo vale si las tres herramientas responden lo mismo. Los resultados se comparan celda a celda contra Pandas como referencia, con tolerancia numérica en las columnas de porcentaje, y se verifica el cuadre aritmético: la suma de viajes de Q1 debe igualar el número de filas del archivo, y las llegadas de Q2 deben igualar el total de Q3. \textbf{Las tres herramientas produjeron resultados idénticos en todas las versiones que completaron.}
